% ============================================================
% PHẦN 3.10 — HOẠT ĐỘNG MARKETING TRONG GIAI ĐOẠN ĐẦU
% ------------------------------------------------------------
% Thuộc: PA3-04
% Người phụ trách chính: Nguyễn Ngọc Cảnh – Operations Lead
% Phối hợp: Nguyễn Anh Thư, Lê Nguyên Nhật
%
% CÂU HỎI CẦN TRẢ LỜI:
%   - Trước ngày ra mắt cần thực hiện hoạt động gì?
%   - Trong tuần ra mắt cần thực hiện hoạt động gì?
%   - Sau ra mắt cần chăm sóc lead và lấy phản hồi như thế nào?
%   - Khi nào nhóm dừng một chiến dịch hoặc thay đổi thông điệp?

% OUTPUT BẮT BUỘC:
%   1. Lịch hoạt động marketing giai đoạn đầu:
%      - Bài giới thiệu vấn đề
%      - Video demo
%      - Tình huống sử dụng (use case)
%      - Workshop
%      - Tiếp cận trực tiếp (outreach)
%      - Mời pilot
%      - Nội dung phản hồi hoặc case study

%   2. Quy trình vận hành pilot:
%      Sàng lọc khách hàng
%        → Tiếp nhận dữ liệu
%        → Chuẩn hóa dữ liệu
%        → Tạo nội dung
%        → Duyệt
%        → Chạy thử
%        → Ghi nhận số liệu
%        → Phỏng vấn sau thử nghiệm
%        → Đề xuất trả phí

% TIÊU CHÍ HOÀN THÀNH:
%   - Không viết "tăng nhận diện" mà không có hành động cụ thể
%   - Mỗi hoạt động phải tạo ra một đầu ra kiểm tra được
%   - Kế hoạch phù hợp với giai đoạn pilot có hỗ trợ, chưa phải SaaS tự phục vụ
% ============================================================

% TODO (Nguyễn Ngọc Cảnh): Viết nội dung mục 3.10
% Lấy kênh đã được chốt từ 3.6 và thông điệp từ 3.5 để tạo lịch nội dung cụ thể


% ============================================================
% PHẦN 3.10 — HOẠT ĐỘNG MARKETING TRONG GIAI ĐOẠN ĐẦU
% ============================================================

\subsection{Hoạt động marketing trong giai đoạn đầu}

Trong giai đoạn đầu (MVP), mục tiêu của hoạt động marketing không phải là phủ sóng diện rộng, mà là tìm kiếm, giáo dục và thuyết phục thành công nhóm khách hàng tiên phong (các trung tâm ngoại ngữ tại TP.HCM quy mô 3-5 nhân sự) tham gia chạy thử nghiệm (Pilot). Các hoạt động được thiết kế gắn liền với đầu ra có thể đo lường được.

\subsubsection{Lịch hoạt động marketing giai đoạn đầu}

\textbf{A. Trước ngày ra mắt (Tuần 1 - Tuần 2: Tạo mồi nhử và vật liệu truyền thông)}
\begin{itemize}
    \item \textbf{Bài giới thiệu vấn đề:} Đăng tải các bài phân tích ngắn trên các cộng đồng quản lý giáo dục/chủ trung tâm ngoại ngữ. Nội dung xoáy thẳng vào nỗi đau: ``Sự lãng phí chi phí nhân sự trực ca đêm và rủi ro rớt lead khi chạy Ads tuyển sinh".
    \begin{itemize}
        \item \textit{Đầu ra kiểm tra được:} 02 bài post trên hội nhóm, đạt tối thiểu 15 bình luận thảo luận từ đúng đối tượng ICP.
    \end{itemize}
    \item \textbf{Video demo:} Quay màn hình thực tế (screen-recording) thể hiện luồng làm việc của 1Stream: từ một bảng giá học phí Excel $\rightarrow$ Video AI giới thiệu khóa học $\rightarrow$ AI Sales Agent tự động trả lời bình luận hỏi giá. Đăng tải lên Landing page và kênh Facebook/TikTok của dự án.
    \begin{itemize}
        \item \textit{Đầu ra kiểm tra được:} 01 Video demo hoàn chỉnh (dưới 2 phút) được nhúng thành công lên Landing page.
    \end{itemize}
    \item \textbf{Tình huống sử dụng (Use case):} Đóng gói một tài liệu ngắn mô phỏng cách một trung tâm tiếng Anh nhỏ dùng 1Stream để treo livestream 30 giờ/tuần mà không cần thuê thêm host.
    \begin{itemize}
        \item \textit{Đầu ra kiểm tra được:} 01 file PDF "Sổ tay tối ưu chi phí livestream tuyển sinh" dùng làm Lead Magnet (quà tặng đổi lấy email/Zalo của khách hàng).
    \end{itemize}
\end{itemize}

\textbf{B. Trong tuần ra mắt (Tuần 3: Tiếp cận và chuyển đổi)}
\begin{itemize}
    \item \textbf{Workshop chuyên đề (Kênh phễu giữa):} Tổ chức một buổi chia sẻ trực tuyến (Zoom/Google Meet) với chủ đề "Tối ưu chi phí livestream và chống rò rỉ Lead ca tối". Thực hiện demo sản phẩm trực tiếp (Live Demo) ngay trong buổi này.
    \begin{itemize}
        \item \textit{Đầu ra kiểm tra được:} Tổ chức thành công 01 buổi workshop, thu hút 15-20 người tham gia thực tế, thu về ít nhất 3 yêu cầu dùng thử.
    \end{itemize}
    \item \textbf{Tiếp cận trực tiếp (Outreach - ABM):} Triển khai kịch bản nhắn tin/gọi điện trực tiếp (Cold Outreach) qua LinkedIn, Zalo hoặc Email tới danh sách 50-80 tài khoản khách hàng mục tiêu đã chấm điểm ICP.
    \begin{itemize}
        \item \textit{Đầu ra kiểm tra được:} Gửi thành công 50 tin nhắn tiếp cận, xác nhận được 5-8 lịch hẹn gọi Demo (Discovery Call) 1-1.
    \end{itemize}
    \item \textbf{Mời pilot:} Đưa ra lời kêu gọi hành động (CTA) rõ ràng cuối buổi Demo/Workshop: Mời khách hàng trải nghiệm gói vận hành 30 giờ-kênh/tuần với chi phí ưu đãi (hoặc đặt cọc cam kết).
    \begin{itemize}
        \item \textit{Đầu ra kiểm tra được:} 02 đến 03 khách hàng ký biên bản đồng ý tham gia Pilot và bàn giao người đầu mối làm việc.
    \end{itemize}
\end{itemize}

\textbf{C. Sau ra mắt (Tuần 4 trở đi: Chăm sóc và Chứng minh giá trị)}
\begin{itemize}
    \item \textbf{Nội dung phản hồi hoặc Case study:} Sau khi phiên pilot đầu tiên kết thúc, tổng hợp số liệu thực tế (số giờ tiết kiệm được, số câu hỏi AI trả lời đúng, số lead thu về) thành một bài học thành công (Case study).
    \begin{itemize}
        \item \textit{Đầu ra kiểm tra được:} 01 bài báo cáo Case Study (ẩn danh hoặc công khai tên trung tâm nếu được cho phép) gửi lại cho toàn bộ danh sách khách hàng đang ở trạng thái "Đang cân nhắc" (Nurturing).
    \end{itemize}
\end{itemize}

\subsubsection{Quy trình vận hành pilot (Mô hình Service-Enabled Software)}

Vì sản phẩm ở giai đoạn này chưa phải là SaaS tự phục vụ hoàn toàn, quy trình Pilot được thiết kế để đội ngũ 1Stream trực tiếp hỗ trợ, giúp giảm rủi ro lỗi kỹ thuật và giảm "chi phí học hỏi" cho khách hàng.

\begin{enumerate}
    \item \textbf{Sàng lọc khách hàng:} Đối chiếu thông tin khách hàng đăng ký với Ma trận chấm điểm ICP (Mục 3.3). Chỉ duyệt các trung tâm đạt >75 điểm (có đủ dữ liệu tĩnh, có nhân sự phối hợp, đúng ngành).
    \begin{itemize}
        \item \textit{Đầu ra:} Bảng chấm điểm duyệt điều kiện tham gia Pilot.
    \end{itemize}
    \item \textbf{Tiếp nhận dữ liệu:} Khách hàng điền thông tin vào "Biểu mẫu Onboarding" do 1Stream cung cấp, bao gồm: bảng học phí, lịch khai giảng, ưu đãi, FAQ và kịch bản tư vấn chuẩn.
    \begin{itemize}
        \item \textit{Đầu ra:} Thư mục lưu trữ dữ liệu thô của khách hàng.
    \end{itemize}
    \item \textbf{Chuẩn hóa dữ liệu:} Đội ngũ Vận hành (Operations) của 1Stream rà soát, làm sạch và chuyển đổi dữ liệu thô thành định dạng Knowledge Base (JSON) để AI có thể hiểu và truy xuất chính xác.
    \begin{itemize}
        \item \textit{Đầu ra:} Bộ cơ sở tri thức (Knowledge Base) đã được chuẩn hóa.
    \end{itemize}
    \item \textbf{Tạo nội dung:} Dùng hệ thống 1Stream để sinh kịch bản livestream, khởi tạo giọng đọc/Avatar AI và ghép video, thêm phụ đề dựa trên dữ liệu đã chuẩn hóa.
    \begin{itemize}
        \item \textit{Đầu ra:} Bản nháp Video AI và Kịch bản chữ.
    \end{itemize}
    \item \textbf{Duyệt:} Gửi bản nháp cho Người duyệt nội dung của trung tâm kiểm tra. Giới hạn chỉnh sửa tối đa 2 lần để đảm bảo tiến độ.
    \begin{itemize}
        \item \textit{Đầu ra:} Email/Tin nhắn xác nhận "Approve" (Đồng ý phát sóng) từ khách hàng.
    \end{itemize}
    \item \textbf{Chạy thử:} Đội ngũ 1Stream thiết lập luồng phát sóng, kết nối API lên kênh Facebook/TikTok của khách hàng và giám sát trạng thái hệ thống trong suốt thời gian diễn ra phiên live.
    \begin{itemize}
        \item \textit{Đầu ra:} Phiên livestream diễn ra thành công trên kênh của khách hàng, không bị gián đoạn tín hiệu.
    \end{itemize}
    \item \textbf{Ghi nhận số liệu:} Xuất nhật ký (log) của hệ thống: Thời lượng phát sóng thực tế, tổng số bình luận, số câu AI trả lời đúng, số trường hợp phải chuyển cho nhân viên thật, số Lead thu được.
    \begin{itemize}
        \item \textit{Đầu ra:} Báo cáo hiệu suất phiên Livestream (Session Report).
    \end{itemize}
    \item \textbf{Phỏng vấn sau thử nghiệm:} Đặt lịch gọi 15 phút với chủ trung tâm hoặc nhân viên trực page để đánh giá độ hài lòng: AI trả lời có tự nhiên không? Áp lực công việc có giảm không?
    \begin{itemize}
        \item \textit{Đầu ra:} Biên bản phỏng vấn / Khảo sát CSAT.
    \end{itemize}
    \item \textbf{Đề xuất trả phí:} Dựa trên báo cáo hiệu suất và phản hồi tích cực, đội ngũ Sales gửi hợp đồng đề xuất chuyển đổi sang Gói dịch vụ trả phí (Starter/Standard).
    \begin{itemize}
        \item \textit{Đầu ra:} Hợp đồng được ký kết, hoặc lý do từ chối rõ ràng được ghi nhận vào CRM.
    \end{itemize}
\end{enumerate}

\subsubsection{Chăm sóc Lead và thu thập phản hồi}

\begin{itemize}
    \item \textbf{Chăm sóc Lead (Lead Nurturing):} Các khách hàng rớt ở khâu "Sàng lọc" (do chưa sẵn sàng dữ liệu) sẽ được đưa vào chuỗi chăm sóc tự động qua Email/Zalo. Nhóm sẽ gửi tặng họ "Bộ Template Excel chuẩn hóa dữ liệu tuyển sinh" để họ tự hoàn thiện nội bộ, tạo tiền đề tiếp cận lại vào tháng sau.
    \item \textbf{Thu thập phản hồi (Feedback Loop):} Áp dụng nguyên tắc phản hồi nhanh gọn (micro-feedback). Thay vì gửi các form khảo sát dài, nhóm vận hành sẽ tương tác trực tiếp qua Zalo group chung với khách hàng ngay sau mỗi ca live: \textit{"Hôm nay AI nhận diện 12 lead hỏi học phí lớp IELTS, bên mình đánh giá câu trả lời đã đúng giọng điệu trung tâm chưa ạ?"}. Mọi góp ý (dù là phàn nàn) đều được ghi nhận vào Ticket hệ thống để CTO điều chỉnh prompt ngay lập tức.
\end{itemize}

\subsubsection{Tiêu chí dừng chiến dịch hoặc thay đổi thông điệp (Pivot/Kill Triggers)}

Để bảo vệ nguồn lực giới hạn, nhóm marketing sẽ liên tục đo lường dựa trên bộ KPI (Mục 3.12) và áp dụng kỷ luật thép để dừng/đổi chiến dịch nếu vi phạm các ngưỡng sau:

\begin{enumerate}
    \item \textbf{Thông điệp không chạm đúng Nỗi đau:} Nếu chiến dịch "Tiếp cận trực tiếp" (Outreach) gửi đi 50 tin nhắn cho các đối tượng đã lọc kỹ (Core ICP) mà \textbf{tỷ lệ phản hồi (Reply rate) dưới 10\%} sau 1 tuần $\rightarrow$ Cần lập tức dừng gửi. Đổi thông điệp từ việc "Nói về tính năng tạo video AI" sang "Đánh thẳng vào số tiền lương ngoài giờ họ đang lãng phí mỗi tháng".
    \item \textbf{Rào cản quy trình quá lớn:} Nếu khách hàng đồng ý xem Demo (tỷ lệ đặt lịch cao) nhưng ở bước \textbf{"Tiếp nhận dữ liệu" có tỷ lệ rớt (Drop-off) lớn hơn 50\%} (khách hàng phàn nàn việc chuẩn bị dữ liệu quá mệt mỏi) $\rightarrow$ Dừng việc bán gói Pilot hiện tại. Hệ thống phải đổi chiến thuật: Thiết kế sẵn các gói Template kịch bản mẫu cho ngành tiếng Anh để khách chỉ cần điền giá tiền, thay vì bắt khách tự viết từ số không.
    \item \textbf{Chi phí thu Lead (CPL) mất kiểm soát:} Nếu sử dụng ngân sách cho các hoạt động đẩy hội thảo (Workshop) nhưng tính toán lại \textbf{CPL cao hơn 2 lần mức cho phép} hoặc tập khách hàng thu về toàn là cá nhân bán hàng nhỏ lẻ không có ngân sách (Sai ICP) $\rightarrow$ Dừng kênh Workshop đại trà, quay lại chiến dịch "Bắn tỉa" (ABM) từng trung tâm cụ thể.
    \item \textbf{Tỷ lệ Chuyển đổi Trả phí bằng 0:} Nếu hoàn tất 3 đợt Pilot thành công, AI hoạt động tốt, khách hàng khen ngợi nhưng \textbf{không có ai đồng ý đề xuất trả phí} $\rightarrow$ Dừng ngay các chiến dịch Marketing mở rộng phễu. Vấn đề nằm ở Cấu trúc Giá (Pricing) hoặc ROI khách hàng nhận được chưa đủ bù đắp chi phí thuê 1Stream. Phải ngồi lại với 3 khách hàng Pilot để tái cấu trúc gói sản phẩm trước khi tìm khách mới.
\end{enumerate}